\chapter{Objetivos y alcance}
\label{ch:objetivos}

\section{Objetivos del TFG}

El objetivo general de este trabajo es construir un sistema que permita consultar
en lenguaje natural el contenido de las actas del Pleno del Ayuntamiento de Bilbao
(2007--2026) y comparar, sobre el mismo corpus, dos arquitecturas de recuperación
de información complementarias: RAG vectorial y GraphRAG.

De este objetivo general se derivan los siguientes objetivos específicos:

\begin{itemize}
	\item Extraer, a partir de los PDF de las actas, las proposiciones individuales
	del Pleno junto con sus metadatos (fecha, grupo proponente, resultado de la
	votación, página de origen).
	\item Construir una ontología OWL y un grafo de conocimiento RDF que modele
	proposiciones, grupos políticos, concejales, temas y votaciones, enriquecido
	mediante un LLM y consultable por SPARQL con razonamiento automático
	(\textit{roll-up} temático).
	\item Construir un censo de concejales del Pleno (con su grupo y periodo de
	actividad) y una capa determinista de listas nominales de voto, para poder
	responder preguntas sobre personas concretas y no solo sobre grupos.
	\item Implementar un sistema de RAG vectorial sobre el mismo corpus, con
	recuperación híbrida (semántica, literal y temática) y re-ranking.
	\item Implementar un sistema GraphRAG que traduzca la pregunta a SPARQL
	mediante un LLM, ejecute la consulta sobre el grafo y narre la respuesta a
	partir de los datos obtenidos.
	\item Dotar a ambos sistemas de citas de fuente verificables --enlace directo
	a la página del PDF de origen de cada dato mostrado.
	\item Validar la calidad e integridad del grafo de conocimiento mediante
	restricciones formales SHACL, como red de seguridad independiente del
	código que lo construye.
	\item Evaluar y comparar ambos sistemas de forma sistemática, con un conjunto
	de preguntas de referencia con respuesta verificada a mano.
	\item Exponer ambos sistemas mediante una interfaz web conversacional.
\end{itemize}

\section{Alcance del sistema desarrollado}

El sistema final integra los dos motores de recuperación descritos --RAG
vectorial y GraphRAG-- detrás de una única interfaz web conversacional
(Chainlit), donde el usuario puede formular preguntas en castellano y recibir una
respuesta narrada con las fuentes citadas. El grafo de conocimiento incluye,
además de las proposiciones y su clasificación temática, el censo de concejales
del Pleno y las listas nominales de voto cuando el acta las recoge de forma
explícita, y se valida mediante un conjunto de \textit{shapes} SHACL
(\texttt{shapes.ttl}) que comprueban la integridad estructural del grafo ya
construido de forma independiente del código que lo genera.

Queda fuera del alcance de esta versión del sistema el enlazado de las
entidades del grafo con identificadores externos de Wikidata mediante
\texttt{owl:sameAs}. Se valoró explícitamente durante el desarrollo y se
descartó para esta entrega: la cobertura de Wikidata sobre los concejales del
Pleno es muy limitada (solo los alcaldes y un pequeño número de cargos que
alcanzaron relevancia autonómica o estatal, frente a los aproximadamente 94
concejales del censo), por lo que no habría resuelto el hueco real de datos
--la falta de afiliación de partido en una parte del censo-- descrito en el
Capítulo~\ref{ch:conclusiones}. El enlazado sí tendría valor claro sobre
entidades no personales mencionadas en las actas (empresas, organizaciones,
lugares), donde Wikidata ofrece buena cobertura y datos adicionales como
coordenadas geográficas; se retoma como mejora futura en el
Capítulo~\ref{ch:conclusiones}.
